\chapter{Модель канала передачи данных}

Реализовать модель многолучевого канала передачи с добавлением аддитивного белого гауссова шума заданной мощности.

\section{Описание модели многолучевого канала передачи}

Модель многолучевого канала передачи включает в себя несколько путей (лучей) распространения сигнала, которые складываются на входе приёмника. Задача модели канала передачи состоит в имитации искажений, возникающих при передаче сигнала в условиях сложной местности (например, городской застройки), где появляется эффект многолучевости.
\begin{figure}[H]
    \centering
    \includegraphics[width=1.0\textwidth]{multipath.png}
    \caption{Схема многолучевого распространения сигнала}
    \label{fig:multipath}
\end{figure}
Количество лучей определим параметром $N_{\text{л}}$. Длину сигнала с выхода передатчика обозначим символом $L$. В свою очередь, каждый луч имеет протяжённость $D$ (в метрах), исходя из которой рассчитываются величины затухания $G$ и задержки $\tau$ для каждого луча. Примем условие, что параметр $D$ выбирается случайным образом в диапазоне от 10 до 500 метров при каждом запуске симуляции с помощью функции рандомайзера.

Луч с минимальным $D$ будет считаться прямым, а его индекс в модели примем $i = 1$. Тогда задержка $\tau_i$ (в дискретных отсчётах) прихода $i$-го луча относительно прямого луча может быть найдена по формуле:

\begin{equation}
\tau_i = \frac{D_i - D_1}{c T_s}, \quad i = 1, 2, \ldots, N_{\text{л}}
\end{equation}

где $c$ — скорость света, а $T_s$ — длительность дискретного отсчёта в модели, которая может быть найдена как $T_s = 1 / B$, где $B$ — полоса информационного сигнала, задаваемая вариантом задания.



После расчёта величина $\tau_i$ должна быть округлена до ближайшего целого значения. В математической модели одним из вариантов реализации задержки в $i$-м луче является добавление целого количества $\tau_i$ нулей перед сигнальными отсчётами:

\begin{equation}
S_i(k) = 
\begin{cases}
0, & \text{если } k \leq \tau_i \\
S_{tx}(k - \tau_i), & \text{если } k > \tau_i
\end{cases}
\end{equation}

где $i = 1, 2, \ldots, N_{\text{л}}$ — номер луча, $k = 1, 2, \ldots, L + \tau_i$ — номер отсчёта в сигнальном векторе, $S_{tx}$ — сигнал с выхода передатчика.

Далее рассчитаем коэффициент ослабления для всех лучей. Для расчёта коэффициента ослабления сигнала в $i$-м луче используется формула ослабления сигнала в вакууме:

\begin{equation}
G_i = \frac{c}{4 \pi D_i f_0}
\end{equation}

где $f_0$ — несущая частота сигнала, задаваемая вариантом задания.

Для внесения ослабления в сигнал требуется умножить весь сигнальный вектор $S_i$ на коэффициент ослабления $G_i$.

После реализации задержки и ослабления сигнала во всех лучах можно получить общий сигнал на выходе канала передачи (на входе приёмника). Для этого требуется просуммировать сигнальные векторы друг с другом. Обозначим сигнал на выходе многолучевого канала переменной $S_{\text{mpy}}$ и рассчитаем его:

\begin{equation}
S_{\text{mpy}} = \sum_{i=1}^{N_{\text{л}}} G_i S_i
\end{equation}

На этом часть модели добавления мультипликативной помехи можно считать завершённой. Далее реализуем добавление шума.

\section{Добавление аддитивного белого гауссова шума}

Аддитивный белый гауссов шум $\mathbf{n}$ накладывается на сигнал $\mathbf{S_{\text{mpy}}}$ с помощью сложения шумового вектора с этим сигналом. Для формирования шумового вектора $\mathbf{n}$ в модели может быть использована встроенная функция. Пример формирования вектора с шумом представлен ниже:

\begin{equation}
\mathbf{n} = \text{wgn}(M, 1, N_0)
\end{equation}

где $M$ — длина вектора с шумом, а $N_0$ — спектральная плотность мощности шума в дБ. Учитывайте, что длина $\mathbf{n}$ должна совпадать с длиной $\mathbf{S_{\text{mpy}}}$.

Для управления мощностью шума в модели используется переменная спектральной плотности мощности шума $N_0$ (дБ), которая задаётся пользователем в качестве входного параметра. Добавление шума на сигнал реализуется с помощью формулы:

\begin{equation}
\mathbf{S_{\text{rx}}} = \mathbf{S_{\text{mpy}}} + \mathbf{n}
\end{equation}

Сигнальный вектор на выходе блока многолучевого канала должен иметь такую же длительность, как и на входе, потому что приёмник обрабатывает сигнал идентичной длительности. Чтобы это осуществить, можно "вырезать" из сигнального массива $\mathbf{S_{\text{rx}}}$ первые $L$ отсчётов и подать их на выход модели канала передачи.

\section{Параметры модели}

В этой модели пользователем задаются параметры, влияющие на тип многолучевого канала:

\begin{itemize}
    \item Количество лучей в многолучевой модели — параметр $N_{\text{В}}$;
    \item Мощность АБГШ в дБ — параметр $N_0$.
\end{itemize}

Для математического расчёта в модели используются следующие константы:

\begin{itemize}
    \item Скорость света $c = 3 \cdot 10^8$ м/с;
    \item Несущая частота (Гц) — задаётся вариантом из таблицы;
    \item Полоса сигнала (Гц) — задаётся вариантом из таблицы.
\end{itemize}

\section{Структурная схема модели}

Структурная схема модели многолучевого канала передачи с АБГШ приведена на рисунке ниже.

\begin{figure}[H]
    \centering
    \includegraphics[width=0.9\textwidth]{channel_model.png}
    \caption{Структурная схема модели многолучевого канала передачи с АБГШ}
    \label{fig:channel_model}
\end{figure}

\section{Описание входных и выходных переменных}

\begin{itemize}
    \item \textbf{Входные переменные:} вектор с комплексными символами $S_{tx}$;
    \item \textbf{Выходные переменные:} вектор с комплексными символами $S_{rx}$;
    \item \textbf{Параметры:}
    \begin{itemize}
        \item $N_{\text{B}}$ — количество лучей;
        \item $N_0$ — мощность АБГШ.
    \end{itemize}
\end{itemize}